\documentclass[11pt]{article}
\usepackage[margin=1in]{geometry}
\usepackage{amsmath,amssymb,amsthm}
\setlength{\parindent}{0pt}\setlength{\parskip}{6pt}

\newtheorem{lemma}{Lemma}
\theoremstyle{remark}
\newtheorem*{remark}{Remark}

\newcommand{\E}{\mathbb{E}}
\newcommand{\Var}{\operatorname{Var}}
\newcommand{\Cov}{\operatorname{Cov}}
\newcommand{\tr}{\operatorname{tr}}

\title{Yi's Lemma and a SURE Form for Ledoit--Wolf Shrinkage}
\author{Jack Wang}
\date{}

\begin{document}
\maketitle
\vspace{-2em}

\section{Motivation}

Under the guidance of Professor Rosenman, the first stage of this project
studied shrinkage in the classical normal-means problem. For the linear rule
$\hat\theta = (1-\lambda)X$, the correction term is $g(X) = -\lambda X$, and
most terms in Stein's unbiased risk estimate (SURE) simplify directly, so the
MSE-optimal amount of shrinkage can be displayed graphically across different
dimensions and directions. The James--Stein rule
$\hat\theta = \bigl(1 - c/\|X\|^2\bigr)X$ admits the same kind of
simplification and can be plotted in the same way.

The central question of the thesis is whether the same program carries over
to covariance estimation: can the Ledoit--Wolf (LW) shrinkage estimator be
written in SURE form? This note derives the key identity, due to Yi and Zou,
and applies it to the LW estimator.

\section{Setup}

Let $\Sigma = (\sigma_{ij})$ be the true $p\times p$ covariance matrix, $S =
(s_{ij})$ the sample covariance matrix, and $\hat\Sigma = (\hat\sigma_{ij})$
any estimator of $\Sigma$. The loss is the squared Frobenius norm. Recall that
for a matrix $A$,
\[
\|A\|_F = \Bigl(\sum_{i,j} |a_{ij}|^2\Bigr)^{1/2}
        = \sqrt{\tr(A^{H}A)}
        = \Bigl(\sum_k \sigma_k^2(A)\Bigr)^{1/2},
\]
where $A^H$ is the conjugate transpose and $\sigma_k(A)$ are the singular
values of $A$. Because the loss is the \emph{squared} norm, the square root
drops out and the risk decomposes entrywise:
\[
R(\hat\Sigma) = \E\|\hat\Sigma - \Sigma\|_F^2
              = \sum_{i,j} \E(\hat\sigma_{ij} - \sigma_{ij})^2 .
\]

\textbf{Assumption.} The sample covariance is unbiased, $\E S = \Sigma$, that is,
$\E s_{ij} = \sigma_{ij}$ for every $(i,j)$.

\begin{remark}
This holds, for example, when the daily observations $x_1,\dots,x_T \in
\mathbb{R}^p$ (one vector per day $t$) are i.i.d.\ across days with mean zero
and $S = \frac1T\sum_t x_t x_t^\top$, since then
$\E s_{ij} = \frac1T\sum_t \E[x_{ti}x_{tj}] = \sigma_{ij}$. If the mean is
estimated, the demeaned estimator needs the divisor $T-1$ instead of $T$ to
stay unbiased. No distributional assumption (such as normality) is needed.
\end{remark}

\section{Yi's lemma}

\begin{lemma}[Yi and Zou]\label{lem:yi}
Under the assumption $\E S = \Sigma$, for any estimator $\hat\Sigma$,
\[
\E\|\hat\Sigma - \Sigma\|_F^2
= \E\|\hat\Sigma - S\|_F^2
  - \sum_{i,j}\Var(s_{ij})
  + 2\sum_{i,j}\Cov(\hat\sigma_{ij}, s_{ij}).
\]
\end{lemma}

\begin{proof}
Fix an entry $(i,j)$ and write $\hat\sigma$, $s$, $\sigma$ for
$\hat\sigma_{ij}$, $s_{ij}$, $\sigma_{ij}$. Adding and subtracting $s$,
\[
\hat\sigma - \sigma = (\hat\sigma - s) + (s - \sigma),
\]
so
\[
(\hat\sigma - \sigma)^2
= (\hat\sigma - s)^2 + (s - \sigma)^2 + 2(\hat\sigma - s)(s - \sigma).
\]
We take expectations term by term.

\emph{First term.} $\E(\hat\sigma - s)^2$ is kept as it is.

\emph{Second term.} Expanding and using $\E s = \sigma$,
\[
\E(s-\sigma)^2 = \E s^2 - 2\sigma\,\E s + \sigma^2
               = \E s^2 - \sigma^2
               = \E s^2 - (\E s)^2
               = \Var(s).
\]

\emph{Cross term.} Distributing,
\[
\E\bigl[(\hat\sigma - s)(s - \sigma)\bigr]
= \E\bigl[\hat\sigma(s-\sigma)\bigr] - \E\bigl[s(s-\sigma)\bigr].
\]
For any random variables $A$ and $B$,
\[
\Cov(A,B) = \E\bigl[(A - \E A)(B - \E B)\bigr]
          = \E\bigl[A(B - \E B)\bigr] - \E A\cdot\E\bigl[B - \E B\bigr]
          = \E\bigl[A(B - \E B)\bigr],
\]
since $\E[B - \E B] = 0$. Taking $B = s$, so that $\E B = \sigma$, and first
$A = \hat\sigma$ and then $A = s$, gives
\[
\E\bigl[\hat\sigma(s-\sigma)\bigr] = \Cov(\hat\sigma, s),
\qquad
\E\bigl[s(s-\sigma)\bigr] = \Cov(s,s) = \Var(s).
\]
Hence the cross term equals $\Cov(\hat\sigma, s) - \Var(s)$.

\emph{Combining.}
\[
\E(\hat\sigma - \sigma)^2
= \E(\hat\sigma - s)^2 + \Var(s) + 2\Cov(\hat\sigma, s) - 2\Var(s)
= \E(\hat\sigma - s)^2 - \Var(s) + 2\Cov(\hat\sigma, s).
\]
Summing over all $(i,j)$, and noting that
$\sum_{i,j}\E(\hat\sigma_{ij} - s_{ij})^2 = \E\|\hat\Sigma - S\|_F^2$,
yields the lemma.
\end{proof}

\begin{remark}
The identity holds for \emph{every} estimator $\hat\Sigma$, whatever function
of the data it is. All three terms are summed over $(i,j)$; the first sum
simply recombines into a Frobenius norm. The estimator enters only through
$\E\|\hat\Sigma - S\|_F^2$ and the covariance term
$\sum_{i,j}\Cov(\hat\sigma_{ij}, s_{ij})$, which plays the role of the
divergence term in Stein's identity.
\end{remark}

\section{Application to Ledoit--Wolf shrinkage}

The LW estimator shrinks $S$ toward a target $F = (f_{ij})$:
\[
\hat\Sigma = \delta F + (1-\delta) S,
\qquad\text{so}\qquad
\hat\sigma_{ij} = \delta f_{ij} + (1-\delta)s_{ij}.
\]
We evaluate the three terms of Lemma~\ref{lem:yi}.

\emph{First term.} Since $\hat\Sigma - S = \delta(F - S)$,
\[
\E\|\hat\Sigma - S\|_F^2 = \delta^2\,\E\|F - S\|_F^2 .
\]

\emph{Covariance term.} By bilinearity of covariance,
\[
\Cov(\hat\sigma_{ij}, s_{ij})
= \delta\Cov(f_{ij}, s_{ij}) + (1-\delta)\Var(s_{ij}).
\]

\emph{Combining.} The variance terms collect as
$-\Var(s_{ij}) + 2(1-\delta)\Var(s_{ij}) = (1-2\delta)\Var(s_{ij})$, so the
risk as a function of the shrinkage intensity is
\[
R(\delta) = \delta^2\,\E\|F - S\|_F^2
          + (1-2\delta)\sum_{i,j}\Var(s_{ij})
          + 2\delta\sum_{i,j}\Cov(f_{ij}, s_{ij}).
\]

\section{The SURE form}

A SURE is any statistic, computable from the data, whose expectation equals
the risk: $\E[\mathrm{SURE}(\delta)] = R(\delta)$. Each term of $R(\delta)$
therefore has to be replaced by an unbiased, computable counterpart. The
first term needs no work, since $\|F - S\|_F^2$ is computed directly from $S$
and $F$ and is unbiased for $\E\|F - S\|_F^2$. Given unbiased estimators
$\widehat{\Var}(s_{ij})$ and $\widehat{\Cov}(f_{ij}, s_{ij})$,
\[
\mathrm{SURE}(\delta) = \delta^2\,\|F - S\|_F^2
  + (1-2\delta)\sum_{i,j}\widehat{\Var}(s_{ij})
  + 2\delta\sum_{i,j}\widehat{\Cov}(f_{ij}, s_{ij}).
\]

\section{Open step: estimating the variance and covariance terms}

The remaining difficulty is to construct $\widehat{\Var}(s_{ij})$ and
$\widehat{\Cov}(f_{ij}, s_{ij})$ in the matrix setting. Everything above rests
on two assumptions: the observations have mean zero, and they are i.i.d.\
across days. Both are essential to the derivation, but neither is fully
realistic for financial returns.

\begin{remark}[Variance term]
With mean-zero i.i.d.\ data, $s_{ij} = \frac1T\sum_t y_t$ is an average of the
i.i.d.\ products $y_t = x_{ti}x_{tj}$, so
$\Var(s_{ij}) = \frac1T\Var(y_t)$. The usual unbiased sample variance of the
$y_t$, divided by $T$, gives an unbiased estimator:
\[
\widehat{\Var}(s_{ij}) = \frac{1}{T(T-1)}\sum_{t=1}^T \bigl(x_{ti}x_{tj} - s_{ij}\bigr)^2 .
\]
\end{remark}

\begin{remark}[Covariance term]
If the target $F$ is fixed (does not depend on the data), then
$\Cov(f_{ij}, s_{ij}) = 0$ and this term vanishes. For data-dependent targets,
such as the scaled identity $F = \nu I$ with $\nu = \tr(S)/p$, the term is
nonzero and must be estimated; this is the step still to be worked out for
general targets.
\end{remark}

\end{document}
